% NAL 1644, batch 9 — Gallica views 161–163, the last three views of the volume.
% Pass: Opus 5.5 (claude-opus-5-5), 2026-09-23 — first pass, unchecked against the views by a human.
%
% Dating. The catalogue gives no date for this volume (an empty string in
% holdings.json); \dating{} is left empty on purpose, and this pass infers
% none. Nothing on these three views is dated.
%
% What the views are. Colour scans rendered near-greyscale, portrait, about
% 3750–3780 × 5600, one page per view. The leaves are of unequal size and
% edges of neighbouring leaves show above or below the page photographed:
% at the top of v. 161 a strip of a taller leaf with a few pale letters and a
% fragment of a figure (not placed, not transcribed); at the foot of v. 163
% the feet of two longer leaves, the first of which carries « Epicycli » and
% a small « c », i.e. the foot of the page of v. 161, transcribed there only.
% Each view was read from the local mirror as four full-width bands at
% 2400 px, with tight crops at 1200–2400 px for doubtful words and for the
% word under the library's stamp on v. 162. Nothing was fetched from Gallica.
% \page{N} is always the view. No view is skipped: all three carry text.
%
% What the three views hold. Latin, one regular upright hand throughout,
% large and widely spaced, on planetary theory: an eccentric « eccentrulus »
% or small circle, an epicycle, the « altitudo media », the « adlateralis »
% « ad parodos terrenas » and « ad parodos solares », the anomalies « ad
% terram » and « ad solem », worked by proportions of sines (« Vt Sinus
% Totus ad Sinum … ita … »). The leaves carry no author's name, and this pass
% attributes the text to no one beyond the volume's title.
%
%   v. 161   (verso, unnumbered; to all appearances the back of the leaf
%            numbered « 78 » on v. 160) CONTINUES A TEXT BEGUN ON VIEW 160:
%            v. 160 ends on the unfinished statement of « Propositio II »
%            (« In ysdem Planetis Ex datis Anomalijs & prostaphæresi
%            adparentia Vna cum Altitudine media & Adlateralibus ad », read
%            here only to settle the join; it is not transcribed in this
%            batch), and v. 161 opens with its end, « parodos terrenas /
%            Inuenire Adlateralem ad parodos Solares ». Then the data (arcs
%            σ, ρ, τ; X, φ), the restriction to Saturn, Jupiter, Mars and
%            Venus, and a construction: A the earth, B the centre of the
%            « eccentrulus », C, D, E, F, G, with G the centre of the
%            epicycle. The page stops there, the foot blank; the proposition
%            is not finished on it, and v. 162 does not continue it.
%   v. 162   (leaf « 79 ») A new piece under the Greek heading « Δεδομένον »:
%            from the mean altitude of Mercury, the « adlaterales » and the
%            anomalies to the earth and to the sun, the distance of Mercury
%            to the sun is found. A « parasceue » by proportions of sines,
%            the places of « protractio » and « contractio », then the
%            altitude corrected « ad Omœopathiam »; the page runs on into
%            v. 163 in mid-proportion.
%   v. 163   (verso, unnumbered) Continues v. 162: four « factæ » (« Ad
%            factam primam … Quartam »), each a proportion of sines on the
%            lettering H, F, G, E, I, N, O, K, M, then « Ac denique Vt
%            Altitudo media Mercurij », where the text stops. A word
%            « emendata » at the right foot, in blacker ink.
%
% Batch edges.
%   - View 161 continues a text begun on view 160 (above): the second half
%     of the statement of « Propositio II ».
%   - View 163 is the last view of the volume. Its text stops on « Ac
%     denique Vt Altitudo media Mercurij », mid-sentence; the word
%     « emendata » at the foot may be a catchword (an inference). The
%     continuation is not in this volume.
%   - Within the batch: v. 161 → v. 162 is a break (a new heading); v. 162 →
%     v. 163 runs on.
%
% Numbers on the leaves.
%   (a) « 79 » in ink at the top right of the recto of v. 162, a stroke like
%       the « 78 » at the top right of v. 160 (seen to settle the join). The
%       versos (v. 161, 163) carry none. This looks like the continuation of
%       the ink series of rectos that batch 2 read as 9–17 on views 22–39;
%       whether it is one series across the volume, and whose, is for the
%       coordinator to reconcile. No \folio{} is written.
%   (b) « 8 » in ink on v. 161, on a curled paper edge in the left margin,
%       not on the page photographed; its leaf and kind are unknown.
%   (c) « c », small, at the right foot of v. 161 (and seen again at the
%       foot of v. 163 on the projecting leaf). Kind unknown; recorded as a
%       mark.
%   (d) The ordinals of the « factæ » (primam … quartam) and of « Propositio
%       II » are content.
%
% Notation. The diagram letters are capitals, set in math ($HFG$); the E is
% written as an epsilon throughout and set E. Arcs of the data are Greek
% minuscules σ, ρ, τ, set $\sigma$, $\rho$, $\tau$; the letter of the
% « adlateralis » on v. 161 is read φ with doubt and set $\varphi$; the mean
% altitude is X. Equality does not occur; proportions are in words (« Vt …
% ad …, Ita … ad … »). Abbreviation « compl » + superscript « ti » kept as
% compl\textsuperscript{ti}. No figure is drawn on these three views; none
% had to be described.
%
% Spelling and hand. Latin as written: u for v and V for U as on the page
% (« Vna », « Vt », « Inuenire », « Interuectus »); « ysdem »; « parodo »;
% « Methatesis »; « Omœopathiam ». The diphthong is an a with a hook, set æ.
% A final i with a descender (« hypothesj », « Epicyclj ») is set i; the
% double ij of « Mercurij », « Anomalijs » is kept. The a of « ad » is often
% open and looks like e; read « ad ». The long s is set s. Line breaks and
% the page's indentation are not kept; each sentence-group of the page is a
% paragraph.

% Coordinator's reconciliation (2026-09-23), after all nine batches:
%   One ink series of leaf numbers runs top right of the rectos through the
%   whole volume, with no gap: 1–8 and « 5 bis » (v. 5–20), 9–17 (v. 22–39),
%   18–27 (v. 41–59), 28–37 (v. 61–79), 38–47 and « 46 bis » (v. 81–100),
%   « 47 bis » and 48–56 (v. 102–120), 57–66 (v. 122–140), 67–78
%   (v. 142–160), 79 (v. 162). It crosses the three pieces of the volume, and
%   the writer cites it himself (« fol. 18. », v. 43; « pag. 48 », v. 100), so
%   it is the writer's or copyist's count of leaves, not the library's. Being
%   ink, it gets no \folio{}, in any batch.
%   The pieces: « Francisci Vietæ De Recognitione Æquationum tractatus »
%   (heading on v. 5), Cap. 1–XXI, v. 5–68; « Francisci Vietæ. De æquationum
%   Emendatione tractatus secundus » (heading on v. 69), Cap. I–XIIII, ending
%   « Finis » on v. 141; then astronomy without a name, v. 142–163 (lunar
%   prostaphaereses, Tycho's lunar hypothesis restated, « Harmonici
%   Ptolemaici Constructio Caput XII », Mercury). Two doubtful notes may mark
%   the treatises as a copy: « transcripta ex exemplari » (v. 100) and
%   « deerant hæc in exemplari » (v. 114). The name on the leaves is in the
%   headings; who wrote the hands is not settled.

\documentclass[11pt,a4paper]{article}
% The preamble shared by every transcription of the mathematicians' manuscripts in Gallica.
%
% It rests on one idea: **uncertainty compounds, it does not resolve.** A
% transcription that smooths over a doubtful word has destroyed the only thing
% it was made for — a reader must be able to tell, at every point, what was on
% the page and what was supplied. The apparatus macros below are therefore the
% critical apparatus, not decoration.
%
% The same commands are recognised by `scripts/render.mjs`, which produces the
% site's reading view, and by `scripts/tei.mjs`. A macro added here must be
% added there too, or rendering fails loudly — which is the wanted behaviour.
%
% Comments here are English, like the rest of the repository. The documents
% this preamble sets are French, and that is deliberate: see the skills.

\ProvidesPackage{germain}[2026/09/15 Transcriptions of mathematicians' manuscripts in Gallica]

\usepackage{amsmath,amssymb,amsthm}
\usepackage{mathrsfs}
% Kept from the parent preamble so the renderer's diagram support stays
% available; these manuscripts draw no commutative diagrams, and nothing here uses it.
\usepackage{tikz-cd}
\usepackage[english,french]{babel}
\usepackage{fontspec}

% --- Greek ------------------------------------------------------------------
% The text face stays fontspec's default, Latin Modern, which has no Greek:
% XeTeX drops every glyph it lacks with a line in the log and nothing on the
% page, and the Greek words the manuscripts quote vanished from the PDFs. So
% the Greek and Greek Extended blocks get a XeTeX character class of their own,
% and entering or leaving a run of it switches the family to CMU Serif — the
% same Computer Modern design, with polytonic Greek — and back. Only the
% family changes: a Greek word in \textit or \textbf keeps its shape and series.
% The transitions to and from every other class carry the tokens that class
% has towards an ordinary letter, so French spacing before « ; » or inside
% « » still holds after a Greek word. No grouping, so no brace or \endgroup
% can come out unbalanced; maths is left alone, since XeTeX inserts nothing
% there. Kept identical in `exercices.sty`.
\makeatletter
\newfontfamily\germainGreekfont{cmun}[
  Extension=.otf, NFSSFamily=germaingreek,
  UprightFont=*rm, ItalicFont=*ti, SlantedFont=*sl,
  BoldFont=*bx, BoldItalicFont=*bi, BoldSlantedFont=*bl]
\newXeTeXintercharclass\germain@greek
\def\germain@greekrange#1#2{%
  \@tempcnta=#1\relax
  \loop
    \XeTeXcharclass\@tempcnta=\germain@greek
    \ifnum\@tempcnta<#2\advance\@tempcnta\@ne\repeat}
\germain@greekrange{"0370}{"03FF}
\germain@greekrange{"1F00}{"1FFF}
\def\germain@greekprev{\rmdefault}
\def\germain@greekin{%
  \edef\germain@greekprev{\f@family}\fontfamily{germaingreek}\selectfont}
\def\germain@greekout{\fontfamily{\germain@greekprev}\selectfont}
\def\germain@greekjoin{%
  \XeTeXinterchartokenstate=\@ne
  \@tempcnta=\z@
  \loop
    \ifnum\@tempcnta=\germain@greek\else
      \XeTeXinterchartoks\germain@greek\@tempcnta=\expandafter{%
        \expandafter\germain@greekout\the\XeTeXinterchartoks\z@\@tempcnta}%
      \XeTeXinterchartoks\@tempcnta\germain@greek=\expandafter{%
        \the\XeTeXinterchartoks\@tempcnta\z@\germain@greekin}%
    \fi
    \ifnum\@tempcnta<4095\advance\@tempcnta\@ne\repeat}
% babel-french sets its own classes, and saves and restores the interchar
% state, each time a language is selected: join again after it does.
\addto\extrasfrench{\germain@greekjoin}
\addto\extrasenglish{\germain@greekjoin}
\AtBeginDocument{\germain@greekjoin}
\makeatother

\usepackage{geometry}
\geometry{a4paper,margin=25mm,marginparwidth=28mm}
\usepackage{marginnote}
\usepackage{xcolor}
\usepackage[normalem]{ulem}
\setlength{\emergencystretch}{3em}
\usepackage{hyperref}
\hypersetup{colorlinks=true,linkcolor=black,urlcolor=black,pdfborder={0 0 0}}

% --- Batch metadata ---------------------------------------------------------
% Taken as they stand by the rendering script, which shows them at the head of
% the reading view. They fix what the file claims to cover. The macro names are
% the parent project's, so that the scripts stay shared; \folder here means the
% volume's slug — `fr-9115` — and \pages counts Gallica views.

\newcommand{\folder}[1]{\gdef\grothFolder{#1}}
\newcommand{\batch}[1]{\gdef\grothBatch{#1}}
\newcommand{\pages}[2]{\gdef\grothFirst{#1}\gdef\grothLast{#2}}
\newcommand{\foldertitle}[1]{\gdef\grothFoldertitle{#1}}

% The holder's own shelfmark, printed as they write it: « Français 9115 ».
\newcommand{\shelfmark}[1]{\gdef\grothShelfmark{#1}}

% The dating, copied verbatim from the holder's catalogue — never inferred
% here. « XIXe siècle » is the BnF's; a pass that dates a leaf from its content
% says so in a \note{}, as its own inference.
\newcommand{\dating}[1]{\gdef\grothDating{#1}}

% The status notice, printed in the title block and repeated as a diagonal
% watermark on every page of the PDF. The manuscripts are in the public
% domain; what this line declares is not a right but a quality — an
% unverified first pass by a machine. The skills write the line; a file
% without it gets no watermark, so leaving it out is visible in the output.
\newcommand{\watermark}[1]{\gdef\grothWatermark{#1}}

\gdef\grothFolder{?}\gdef\grothBatch{}\gdef\grothFirst{?}\gdef\grothLast{?}%
\gdef\grothFoldertitle{}\gdef\grothShelfmark{}\gdef\grothDating{}\gdef\grothWatermark{}

% --- The résumé -------------------------------------------------------------
\newenvironment{resume}
  {\small\setlength{\parskip}{0.45em}}
  {\par\medskip\hrule\medskip}

% --- Keywords ---------------------------------------------------------------
% In English, closing the modernised reading's résumé: the modern vocabulary
% under which to search for what the leaves construct. The manifest extracts
% them to tag the volume — this line is the single source of the tags.
\newcommand{\keywords}[1]{\par\smallskip\noindent
  {\footnotesize\textsc{Keywords} --- \textit{#1}}\par}

% --- The critical apparatus -------------------------------------------------

% The Gallica view number, in the margin. This is the anchor that lets a
% disagreement over a formula be settled by looking at *one* image rather than
% a volume of 750. The site uses it to turn the facsimile as the transcription
% is scrolled. A view is one scanned image: usually one side of a leaf.
\newcommand{\page}[1]{%
  \marginnote{\footnotesize\color{gray!70}\textsf{v.\,#1}}%
  \label{page:#1}\ignorespaces}

% The BnF's pencilled foliation, where the leaf carries it and a pass can read
% it: « 198r ». This is what the literature cites — « ff. 198r–208v » — and
% the one line that turns a citation into a view. Recorded, never inferred:
% a folio a pass cannot read is not written.
\newcommand{\folio}[1]{%
  \marginnote{\footnotesize\color{gray!50}\textsf{f.\,#1}}\ignorespaces}

% The modernised reading's coarser counterpart to \page{}: which views a
% section is drawn from.
\newcommand{\pagerange}[2]{%
  \marginnote{\footnotesize\color{gray!70}\textsf{v.\,#1--#2}}%
  \ignorespaces}

% Illegible: never guessed.
\newcommand{\ill}{\textcolor{red!60!black}{[\,\ldots\,]}}

% A doubtful reading: the word is given, and flagged as uncertain.
\newcommand{\uncertain}[1]{\uline{#1}}

% An editorial addition — a missing letter, an implied word.
\newcommand{\add}[1]{\textcolor{blue!50!black}{[#1]}}

% Struck out by the writer. What was crossed out is part of the document.
\newcommand{\struck}[1]{\sout{#1}}

% The transcriber's note, on the state of the page or the mathematics.
\newcommand{\note}[1]{\par\smallskip{\small\color{gray!60!black}\textit{[#1]}}\par\smallskip}

% A marginal note of hers, distinct from the body of the text.
\newcommand{\marginal}[1]{\par\smallskip\noindent\hspace{1em}%
  {\small\color{gray!60!black}#1}\par\smallskip}

% --- Title ------------------------------------------------------------------
% The title block is French, like the documents themselves.

\AddToHook{shipout/background}{%
  \ifx\grothWatermark\empty\else
    \begin{tikzpicture}[remember picture, overlay]
      \node[rotate=52, scale=3.4, align=center, text=gray!18,
            font=\sffamily\bfseries]
        at (current page.center) {\grothWatermark};
    \end{tikzpicture}%
  \fi}

\AtBeginDocument{%
  \begin{center}
    {\large\bfseries Manuscrits de mathématiciens (Gallica) ---
      \ifx\grothShelfmark\empty\grothFolder\else\grothShelfmark\fi}\\[2pt]
    {\itshape\grothFoldertitle}\\[4pt]
    {\small vues \grothFirst--\grothLast
      \ifx\grothBatch\empty\else\ (lot \grothBatch)\fi}\\[2pt]
    \ifx\grothDating\empty\else
      {\small\color{gray!60!black}Datation du catalogue : \grothDating}\\[2pt]
    \fi
    \ifx\grothWatermark\empty\else
      {\small\bfseries\color{red!45!gray}\grothWatermark}\\[2pt]
    \fi
    {\footnotesize\color{gray!60!black}Transcription établie d'après les images IIIF de Gallica.
     Source gallica.bnf.fr / Bibliothèque nationale de France.\\
     Les lectures incertaines sont soulignées, les illisibles notées [\,\ldots\,],
     les ajouts entre crochets ; \textsf{v.} numérote les vues de Gallica,
     \textsf{f.} la foliotation au crayon de la BnF.}
  \end{center}
  \vspace{1em}}

\endinput


\folder{nal-1644}
\shelfmark{NAL 1644}
\batch{9}
\pages{161}{163}
\dating{}
\watermark{Lecture automatique\\non vérifiée}
\foldertitle{Viète (François). Papiers divers. NAL 1644}

\begin{document}

\page{161}
\note{Verso, sans chiffre. Le feuillet photographié a le bord supérieur déchiré ; au-dessus, dans la bande où dépasse un feuillet plus haut placé dessous, quelques lettres pâles et le fragment d'une figure (un petit cercle pointé au centre, des lettres « a », « A » et d'autres) qui appartiennent à une page que ce lot ne photographie pas ; elles ne sont pas transcrites. Dans la marge de gauche, sur un bord de papier replié, un « 8 » à l'encre, dont le transcripteur ne sait pas à quel feuillet il appartient. Au pied, contre le bord inférieur et à droite, un petit « c » isolé, sans renvoi repéré ; sous ce bord dépasse le pied d'un autre feuillet, sans texte lisible. Les fins de lignes, à droite, touchent le bord du feuillet (« Centru », « normali », « centru ») : les dernières lettres sont cachées ou rognées. Même main, même écriture régulière que la vue 162. La page continue le texte de la vue 160 (feuillet « 78 », recto) : celle-ci s'achève sur l'énoncé, laissé en suspens, de la « Propositio II » (« In ysdem Planetis Ex datis Anomalijs \& prostaphæresi adparentia Vna cum Altitudine media \& Adlateralibus ad »), et la vue 161 en donne la fin, « parodos terrenas / Inuenire Adlateralem ad parodos Solares ». Cette vue est donc, selon toute apparence, le verso du feuillet « 78 ».}

parodos terrenas

Inuenire Adlateralem ad parodos Solares

Sit data Anomalia ad terram Peripheria $\sigma$ ad Solem Peripheria $\rho$ prostaphæresis adparentia peripheria $\tau$
\note{Les trois arcs sont désignés par des lettres grecques minuscules : $\sigma$, $\rho$ (tracé comme un p latin à longue hampe), $\tau$. « prostaphæresis » : la page écrit « prostapharesis », l'æ étant un a à crochet.}

Qualium autem altitudo media $X$ talium data sit adlateralis transversa ad parodos terrenas $\varphi$ \& ideo recta $\varphi$ bis
\note{La lettre de l'« adlateralis » est une boucle haute sur une longue hampe, lue $\varphi$ ; elle revient trois fois sur la page, toujours de même forme, et le transcripteur n'en est pas sûr (une forme de $\varphi$ ou de $\phi$). « ad parodos » : l'a de « ad » est ouvert et se lit presque « ed », ici comme ailleurs sur ces pages ; lu « ad ».}

Oportet in ysdem planetis de quibus \uncertain{emecata} est superior propositio id est Saturno Ioue \& Marte ac Venere Inuenire Adlateralem ad parodos Solares
\note{« emecata » : lecture des lettres (e, m, e, c, a, t, a), qui ne donnent pas de mot latin connu du transcripteur ; le sens attend un participe comme « enucleata » ou « enunciata », que la page ne porte pas. Entre « Superior » et « propositio », une tache d'encre pâle, sans lettre lisible.}

Ex hypothesi igitur qualium $X$ semidiameter homocentri Talium erit $\varphi$ semidiameter Circelli ad Ellipsim seu \uncertain{Eccentruli} Quæritur autem Semidiameter Epicycli
\note{« hypothesi » et « Epicycli » finissent par un i long (« hypothesj », « Epicyclj »), gardé ici comme i. Le mot après « seu » se lit aussi « Eccentrali » : l'a et l'u se confondent dans cette main ; « Eccentruli » est le génitif du « Eccentrulus » écrit quatre lignes plus bas.}

Quocirca esto $A$ terra homocentrus $BZ$ $B$ Centru\add{m} Eccentruli Erit igitur $AB$ linea mediæ \uncertain{Epochæ} planetæ
\note{« Centru » touche le bord du feuillet ; la dernière lettre n'est pas visible. « mediæ », « planetæ » : a final à crochet. « Epochæ » : le mot est écrit « Epocher » avec un point, où le transcripteur lit la finale æ sans en être sûr.}

Describatur Eccentrulus \uncertain{Interuectus} a recta $AB$ In Signis $C$ $D$ illo altiore hoc depressiore Itaque numeretur in antecedentia peripheria $CE$ data $\sigma$ æqualis \& cadat in $CD$ normali\add{s} $EF$ quæ producatur extra circulum in $G$ ut sit $FG$ dupla ipsius $EF$ Erit igitur $G$ centru\add{m} Epicycli
\note{« Interuectus » : la syllabe du milieu est serrée et pourrait se lire « Intersectus » ; lecture douteuse. « normali » et « centru » touchent le bord du feuillet. « æqualis » est écrit « aqualis », avec l'a à crochet. La page s'arrête ici ; le bas du feuillet est blanc. Le texte ne passe pas à la vue 162, qui commence un nouvel énoncé sous un titre grec.}

\page{162}
\note{Recto. En haut à droite, « 79 » à l'encre, de même tracé que le « 78 » de la vue 160 (voir l'en-tête). Le feuillet est le dernier de la pile : sous son bord inférieur, déchiré, le fond est vide, et rien ne dépasse d'un autre feuillet. L'encre du verso (vue 163) transparaît à l'envers sur toute la page ; ces traces ne sont pas transcrites. Le cachet « Bibliothèque nationale, Manuscrits, R.F. » est posé sur la ligne « Ita Adlateralis … ad parodos terrenas » et couvre un mot. Un rectangle gris pâle, à droite, à la hauteur de « Ad protractionem … », ne porte pas d'écriture. Même main que les vues 161 et 163. La page ouvre un nouvel énoncé sous un titre grec ; elle ne continue pas la vue 161, qui s'arrêtait sur la construction du centre de l'épicycle.}

\section*{Δεδομένον}

\note{Titre centré, en grec : « Δεδομένον », « donné » (au sens des \emph{Data}).}

Ex datis Altitudine media Mercurij \& Adlateralibus Vna cum Anomalijs ad Terram \& Solem Inuenitur Distantia Mercurij ad Solem.
\note{« Distantia » : les lettres « -st- » sont sous une tache. Le second « Mercurij » est traversé d'un pâté d'encre, qui laisse lire le mot. Le u et le v sont écrits comme sur la page (« Vna », « Inuenitur », « Vt ») ; il en va de même dans tout le lot.}

Enimuero ad Parasceuem. Erit Vt Sinus Totus

Ad Sinum complementi Anomaliæ duplæ ad Terram

Ita Adlateralis \uncertain{subdupla} ad parodos terrenas.
\note{« subdupla » : la première syllabe est sous le cachet, qui laisse voir la haste d'un s long et la panse d'un b ; « -dupla » est net. La même phrase, en tête de la vue 163, porte « subdupla » en clair. « Adlateralis » : l's final est au bord du cachet.}

Ad protractionem contractionemue Altitudinis Mediæ Vt prodeat Altitudo media æquata
\note{« contractionemue » est écrit en deux morceaux, « contractionem ue ». « Mediæ », « Anomaliæ », « duplæ » et « æquata » : a à crochet, lu æ.}

Erit ut Sinus Totus ad Sinum compl\textsuperscript{ti} $CE$ Ita $BE$ ad $BF$
\note{« compl » suivi d'un t surmonté d'un crochet : abréviation de « complementi ». Les lettres $C$, $E$, $B$, $F$ renvoient à la figure de la vue 160 (feuillet « 78 »), qui porte $B$, $C$, $E$, $F$, $G$ (voir la note de la vue 163) ; l'$E$ est tracé en forme d'epsilon, comme partout dans ce lot.}

Est autem locus protractionis in primo \& ultimo quadrante Anomaliæ duplæ ad Terram Contractionis In tertio \& quarto
\note{Écrit tel quel : le « dernier » quadrant de la protraction est aussi le « quatrième » de la contraction ; on attendrait « secundo \& tertio » pour la contraction. Le transcripteur ne corrige pas. « in primo » : l'i est ouvert et se lit presque « an ».}

Vt ex Diagrammatis per quascumque Methatesis prolatis euidens fit
\note{« Methatesis » : ainsi écrit, pour \emph{metathesis}.}

Deinde emendata ad Omœopathiam altitudine Erit Vt Sinus Totus

Ad Sinum Anomaliæ ad Terram
\note{« Omœopathiam » : l'o et l'e liés, ou « Omæopathiam » ; la ligature est gardée comme œ. La page s'arrête sur « Ad Sinum Anomaliæ ad Terram », sans le troisième terme de la proportion : la phrase continue en tête de la vue 163 (« Ita Adlateralis subdupla ad parodos terrenas »). Le bas du feuillet est blanc.}

\page{163}
\note{Verso, sans chiffre : selon toute apparence le verso du feuillet « 79 » (vue 162), dont la page reprend la phrase. Même main. Le cachet « Bibliothèque nationale, Manuscrits, R.F. » est posé à droite, sur la fin de la ligne « Ita IM ad OM » ; l'encre du recto transparaît faiblement. Sous le bord inférieur du feuillet dépassent les pieds de deux feuillets plus longs : sur le premier, à gauche, le mot « Epicycli », qui est la dernière ligne de la page de la vue 161, où il est transcrit, et à droite le petit « c » du pied de cette même page ; le second est blanc. Plus haut, au pied de la page elle-même, à droite, au-dessous d'un pli horizontal et contre le bord, le mot « emendata », d'une encre plus noire (voir la dernière note de la vue).}

Ita Adlateralis subdupla ad parodo terrenas
\note{« parodo » : ainsi écrit, sans s, là où la vue 162 porte « ad parodos terrenas ». La ligne achève la proportion laissée ouverte au bas de la vue 162 (« Erit Vt Sinus Totus / Ad Sinum Anomaliæ ad Terram »).}

Ad factam primam.

Erit Vt Sinus Totus ad Sinum anguli $HFG$ Ita $EG$ dimidia ad dimidiam $HG$ id est ad $HI$ seu $NO$
\note{Le F de « $HFG$ » porte un trait de plus, comme s'il était récrit sur une autre lettre (un E ?). Le N de « $NO$ » est surchargé d'un point. L'$E$ est tracé en epsilon, comme dans tout le lot.}

Et rursus Vt Sinus Totus

Ad Sinum complementi Anomaliæ ad Terram

Ita Adlateralis Tota

Ad factam Secundam

Erit Vt Sinus Totus ad Sinum Anguli $HGF$ Ita $GF$ ad $FH$

Et Vt Sinus Totus

Ad Sinum Anomaliæ ad Solem.

Ita Adlateralis ad parodos Solares.

Ad factam tertiam

Erit Vt Sinus Totus ad Sinum peripheriæ $KM$ seu anguli $IMO$ Ita $IM$ ad $OM$
\note{« peripheriæ » : le début du mot est récrit sur un premier tracé. « $OM$ » touche le bord du cachet ; les deux lettres se lisent.}

Et rursus Vt Sinus Totus

Ad Sinum Complementi Anomaliæ ad Solem

Ita Adlateralis ad parodos Solares.

Ad factam Quartam

Erit Vt Sinus Totus ad Sinum compl\textsuperscript{ti} peripheriæ $MH$ seu anguli $IMO$ Ita $IM$ ad $IO$
\note{« compl » avec un t surmonté d'un crochet, comme à la vue 162 : « complementi ». Entre le M et le H de « $MH$ », un point. Les « factæ » troisième et quatrième nomment toutes deux l'angle $IMO$, l'une pour le sinus de $KM$, l'autre pour le sinus du complément de $MH$ ; c'est écrit ainsi, et le transcripteur ne corrige pas. Les lettres $H$, $F$, $G$, $E$, $I$, $K$ figurent sur la figure de la vue 160 (feuillet « 78 ») telle que la montre une vignette ; $M$, $N$, $O$ n'y ont pas été vues. Aucune figure n'est dessinée sur les vues 161–163.}

Ac denique Vt Altitudo media Mercurij
\note{La page s'arrête sur ces mots, sans la suite de la proportion ; le reste de la page est blanc jusqu'au pied, où, à droite, contre le bord et sous un pli horizontal, est écrit « emendata », d'une encre plus noire. Le mot peut être une réclame (le premier mot de la page suivante), ou la fin d'une ligne cachée : c'est une inférence. La vue 163 est la dernière du volume : la suite du texte n'y est pas.}

\marginal{emendata}

\end{document}
